\backmatter
\KOMAoptions{bibliography=totoc}
\begin{thebibliography}{34}

\bibitem{src:Birch-and-Swinnerton-Dyer-1965}
B. J. Birch, H. P. F. Swinnerton-Dyer. \emph{Notes on elliptic curves. II}. Journal f\"ur die reine und angewandte Mathematik 218, 79-108. 1965.

\bibitem{src:Bombieri}
E. Bombieri. \emph{Problems of the Millennium: the Riemann Hypothesis}. Clay Mathematics Institute, Millennium Prize Problem statement, \texttt{https:/\allowbreak{}/\allowbreak{}www.claymath.org/\allowbreak{}wp-content/\allowbreak{}uploads/\allowbreak{}2022/\allowbreak{}05/\allowbreak{}riemann.pdf}.

\bibitem{src:Caupin-and-Herbert-2006}
F. Caupin, E. Herbert. \emph{Cavitation in water: a review}. Comptes Rendus Physique 7(9-10), 1000-1017, doi:10.1016/j.crhy.2006.10.015. 2006.

\bibitem{src:Clay-Mathematics-Institute}
Clay Mathematics Institute. \texttt{https:/\allowbreak{}/\allowbreak{}www.claymath.org/\allowbreak{}millennium-problems/\allowbreak{}}.

\bibitem{src:Cook}
Stephen Cook. \emph{The P versus NP Problem}. Clay Mathematics Institute, Millennium Prize Problem statement, \texttt{https:/\allowbreak{}/\allowbreak{}www.claymath.org/\allowbreak{}wp-content/\allowbreak{}uploads/\allowbreak{}2022/\allowbreak{}06/\allowbreak{}pvsnp.pdf}.

\bibitem{src:Cremona-1997}
J. E. Cremona. \emph{Algorithms for Modular Elliptic Curves}. Second edition, Cambridge University Press. 1997.

\bibitem{src:Deligne}
Pierre Deligne. \emph{The Hodge Conjecture}. Clay Mathematics Institute, Millennium Prize Problem statement, \texttt{https:/\allowbreak{}/\allowbreak{}www.claymath.org/\allowbreak{}wp-content/\allowbreak{}uploads/\allowbreak{}2022/\allowbreak{}06/\allowbreak{}hodge.pdf}.

\bibitem{src:Dokchitser-2004}
T. Dokchitser. \emph{Computing special values of motivic L-functions}. Experimental Mathematics 13(2), 137-150. 2004.

\bibitem{src:Duraiswami-2026}
R. Duraiswami. \emph{Self-similar swirl between contracting porous walls: the GD1998 exact Navier--Stokes solution revisited in the similarity variables of the OpenAI 2026 forced blow-up construction}. arXiv:2609.17642v1. 2026.

\bibitem{src:Fefferman}
Charles L. Fefferman. \emph{Existence and Smoothness of the Navier-Stokes Equation}. Clay Mathematics Institute, Millennium Prize Problem statement, \texttt{https:/\allowbreak{}/\allowbreak{}www.claymath.org/\allowbreak{}wp-content/\allowbreak{}uploads/\allowbreak{}2022/\allowbreak{}06/\allowbreak{}navierstokes.pdf}.

\bibitem{src:Formal-Conjectures}
Google DeepMind Formal Conjectures authors. \emph{Lean formalization of the Navier-Stokes existence and smoothness problem statement}. \texttt{https:/\allowbreak{}/\allowbreak{}github.com/\allowbreak{}google-deepmind/\allowbreak{}formal-conjectures/\allowbreak{}blob/\allowbreak{}main/\allowbreak{}FormalConjectures/\allowbreak{}Millenium/\allowbreak{}NavierStokes.lean}.

\bibitem{src:Gibbs-1884}
J. W. Gibbs. \emph{Elements of Vector Analysis, arranged for the use of students in physics}. New Haven, printed 1881--1884, not published. 1884.

\bibitem{src:Gourdon-2004}
X. Gourdon. \emph{The $10^{13}$ first zeros of the Riemann Zeta function, and zeros computation at very large height}. 2004.

\bibitem{src:IAPWS-R12-08}
International Association for the Properties of Water and Steam. \emph{Release on the IAPWS Formulation 2008 for the Viscosity of Ordinary Water Substance}. IAPWS R12-08, Berlin. 2008.

\bibitem{src:IAPWS-R6-95-2018}
International Association for the Properties of Water and Steam. \emph{Revised Release on the IAPWS Formulation 1995 for the Thermodynamic Properties of Ordinary Water Substance for General and Scientific Use}. IAPWS R6-95(2018), Prague. 2018.

\bibitem{src:Jaffe-and-Witten}
Arthur Jaffe, Edward Witten. \emph{Quantum Yang-Mills Theory}. Clay Mathematics Institute, Millennium Prize Problem statement, \texttt{https:/\allowbreak{}/\allowbreak{}www.claymath.org/\allowbreak{}wp-content/\allowbreak{}uploads/\allowbreak{}2022/\allowbreak{}06/\allowbreak{}yangmills.pdf}.

\bibitem{src:Ladyzhenskaya-1967}
O. A. Ladyzhenskaya. \emph{New equations for the description of the motions of viscous incompressible fluids, and global solvability for their boundary value problems}. Trudy Mat. Inst. Steklov. 102, 85--104. 1967.

\bibitem{src:Ladyzhenskaya-1968}
O. A. Ladyzhenskaya. \emph{On modifications of the Navier--Stokes equations for large gradients of the velocities}. Zap. Nauchn. Sem. LOMI 7, 126--154. 1968.

\bibitem{src:Ladyzhenskaya-1969}
O. A. Ladyzhenskaya. \emph{The Mathematical Theory of Viscous Incompressible Flow}. Second English edition, translated by R. A. Silverman and J. Chu, Gordon and Breach. 1969.

\bibitem{src:Leray-1934}
J. Leray. \emph{Sur le mouvement d'un liquide visqueux emplissant l'espace}. Acta Mathematica 63, 193--248, doi:10.1007/BF02547354; translation by R. E. Terrell, arXiv:1604.02484. 1934.

\bibitem{src:Lienstromberg-Schiffer-Schubert-2025}
C. Lienstromberg, S. Schiffer, R. Schubert. \emph{A variational approach to the Navier--Stokes equations with shear-dependent viscosity}. arXiv:2312.03546v2. 2025.

\bibitem{src:Navier-1827}
C.-L. Navier. \emph{M\'emoire sur les lois du mouvement des fluides}. M\'emoires de l'Acad\'emie royale des sciences de l'Institut de France 6, 389--440, read 18 March 1822. 1827.

\bibitem{src:Odlyzko-2001}
A. M. Odlyzko. \emph{The $10^{22}$-nd zero of the Riemann zeta function}. In Dynamical, Spectral, and Arithmetic Zeta Functions, Contemporary Mathematics 290, American Mathematical Society, 139-144. 2001.

\bibitem{src:OpenAI-2026-Euler}
OpenAI. \emph{Finite time blowup for the Euler equation}. \texttt{https:/\allowbreak{}/\allowbreak{}cdn.openai.com/\allowbreak{}pdf/\allowbreak{}315b36cd-ec98-4023-8342-93345194ece1/\allowbreak{}euler.pdf}. 2026.

\bibitem{src:OpenAI-2026-Lean}
OpenAI. \emph{NavierStokesAndEuler}, Lean~4 formalizations of the Navier-Stokes and Euler blowup results. \texttt{https:/\allowbreak{}/\allowbreak{}github.com/\allowbreak{}openai/\allowbreak{}NavierStokesAndEuler}. 2026.

\bibitem{src:OpenAI-2026-Navier-Stokes}
OpenAI. \emph{Finite time blowup for Navier-Stokes}. \texttt{https:/\allowbreak{}/\allowbreak{}cdn.openai.com/\allowbreak{}pdf/\allowbreak{}32d9f210-8b73-45e0-91bc-82a30aef8a9a/\allowbreak{}navier-stokes.pdf}. 2026.

\bibitem{src:Platt-and-Trudgian-2021}
D. Platt, T. Trudgian. \emph{The Riemann hypothesis is true up to $3\cdot 10^{12}$}. Bulletin of the London Mathematical Society 53(3), 792-797. 2021.

\bibitem{src:Quigg-crystallography}
D. Quigg. \emph{Crystallography}. orior research papers, \texttt{https:/\allowbreak{}/\allowbreak{}github.com/\allowbreak{}dstroy0/\allowbreak{}orior}. 2026.

\bibitem{src:Quigg-instruments}
D. Quigg. \emph{The Instruments}. orior research papers, \texttt{https:/\allowbreak{}/\allowbreak{}github.com/\allowbreak{}dstroy0/\allowbreak{}orior}. 2026.

\bibitem{src:Riemann-1859}
Bernhard Riemann. \emph{Handwritten working manuscript, Nachlass Cod. Ms. Riemann 3, Rechenseiten}. Niedersaechsische Staats- und Universitaetsbibliothek Goettingen, Cod. Ms. Riemann 3; scan page 00000038.jpg. 1859.

\bibitem{src:RITS-2026}
NYU Shanghai RITS. \emph{OpenAI Claims a Navier-Stokes Proof, Amid a Dispute Over Credit}. \texttt{https:/\allowbreak{}/\allowbreak{}rits.shanghai.nyu.edu/\allowbreak{}ai/\allowbreak{}openai-navier-stokes-proof-credit-dispute/}. 2026.

\bibitem{src:Stokes-1845}
G. G. Stokes. \emph{On the theories of the internal friction of fluids in motion, and of the equilibrium and motion of elastic solids}. Transactions of the Cambridge Philosophical Society 8, 287--319; reprinted in Mathematical and Physical Papers, vol. I, Cambridge, 1880, 75--129. 1845.

\bibitem{src:Wiles}
Andrew Wiles. \emph{The Birch and Swinnerton-Dyer Conjecture}. Clay Mathematics Institute, Millennium Prize Problem statement, \texttt{https:/\allowbreak{}/\allowbreak{}www.claymath.org/\allowbreak{}wp-content/\allowbreak{}uploads/\allowbreak{}2022/\allowbreak{}05/\allowbreak{}birchswin.pdf}.

\bibitem{src:Wilkins-1998}
Bernhard Riemann, translated by David R. Wilkins. \emph{On the Number of Prime Numbers less than a Given Quantity}. Preliminary version, December 1998, of Ueber die Anzahl der Primzahlen unter einer gegebenen Grosse, Monatsberichte der Berliner Akademie, November 1859.

\end{thebibliography}
